% Методичка требует ответы на контрольные вопросы в отчёте, поэтому раздел включён.
\section{Ответы на контрольные вопросы}
\textbf{1. Почему наивный поиск на случайном тексте работает почти линейно, хотя его оценка
в худшем случае — $O(n \cdot m)$?}

На случайном тексте несовпадение обычно находится на первом-втором символе образца, поэтому
на каждую позицию приходится в среднем меньше двух сравнений, а не $m$. В
таблице~\ref{tab:comparisons} наивный алгоритм на случайном тексте делает 21~084 сравнения
при $n = 16~000$ — того же порядка, что и длина текста.

\textbf{2. Что хранит префикс-функция и почему КМП никогда не возвращается назад по тексту?}

Для каждой позиции образца префикс-функция хранит длину наибольшего собственного префикса,
который совпадает с суффиксом уже прочитанной части. При несовпадении эта длина показывает,
какая часть образца уже совпала с текстом после сдвига, поэтому продолжать сравнение можно с
того же символа текста.

\textbf{3. На каких входах КМП заметно выигрывает у наивного алгоритма, а на каких — нет?}

КМП выигрывает на текстах с длинными частичными совпадениями, например на худшем входе этой
работы. На случайном тексте выигрыша почти нет: при $n = 16~000$ время в
таблице~\ref{tab:time} — 1,437 и 1,348~мс.
